\section{Teor\'ia de curvas}
\label{sec:teoria_curvas}

\begin{definicion}[Curva]
  \label{def:curva}
  Una \textbf{curva} en $\R^n$ es una funci\'on
  $\vb{r}:I\subset\R\rightarrow\R^n$. Si $I=[a,b]$ con $a<b$, entonces
  $\vb{r}(a)$ y $\vb{r}(b)$ son los extremos de la curva.
\end{definicion}

\begin{ejemplo}[Recta]
  \label{ej:recta}
  \begin{equation}
    \label{eq:recta}
    \boldsymbol{\ell}(t)=\vb{r}_0+\vh{u}\,t,
    \qquad t\in\R .
  \end{equation}
\end{ejemplo}

\begin{ejemplo}[Elipse]
  \label{ej:elipse}
  Con $I=[0,2\pi]$,
  \begin{equation}
    \label{eq:elipse}
    \vb{r}(\vph)=
    \begin{pmatrix}
      a\cos\vph\\
      b\sin\vph
    \end{pmatrix}.
  \end{equation}
\end{ejemplo}

\begin{ejemplo}[H\'elice]
  \label{ej:helice}
  \begin{equation}
    \label{eq:helice}
    \vb{r}(t)=
    \begin{pmatrix}
      R\cos(\om t)\\
      R\sin(\om t)\\
      a t
    \end{pmatrix},
    \qquad t\in\R .
  \end{equation}
\end{ejemplo}

\begin{observacion}
  Una parametrizaci\'on arbitraria no siempre genera una curva
  \emph{decente}. Cuando la parametrizaci\'on s\'i produce una curva regular
  hablamos de un \textbf{mapeo}.
\end{observacion}

\subsection{Vector tangente}
\label{subsec:vector_tangente}

Definimos el vector tangente a $\vb{r}(t)$ como la velocidad instant\'anea
\begin{equation}
  \label{eq:velocidad}
  \vb{V}=\lim_{\De t\rightarrow 0}
    \frac{\vb{r}(t+\De t)-\vb{r}(t)}{\De t}=\dv{}{t}\vb{r}(t),
\end{equation}
y la rapidez instant\'anea como su norma,
\begin{equation}
  \label{eq:rapidez}
  V(t)=\norm{\vb{V}}=\norm{\dv{}{t}\vb{r}(t)} .
\end{equation}
Si $\vb{V}(t)$ es diferenciable, la aceleraci\'on es
\begin{equation}
  \label{eq:aceleracion}
  \vb{a}=\dv{}{t}\vb{V}=\ddv{}{t}\vb{r}(t).
\end{equation}

\subsection{Longitud de arco}
\label{subsec:longitud_arco}

Podemos aproximar la longitud de una curva por rectas que unen puntos
sucesivos. Partiendo $t\in[t_0,t_f]$ en $N$ trozos
$t_0,t_1,\dots,t_f$ con $\De t_i=t_{i+1}-t_i$,
\begin{equation}
  \label{eq:longitud_aproximada}
  S\approx\sum_{i=1}^{N}\norm{\vb{r}(t_i+\De t_i)-\vb{r}(t_i)} .
\end{equation}
Tomando $N\rightarrow\infty$, equivalentemente $\De t_i\rightarrow 0$,
\begin{equation}
  \label{eq:longitud_limite}
  S=\lim_{N\rightarrow\infty}\sum_{i=1}^{N}
    \norm{\frac{\vb{r}(t_i+\De t_i)-\vb{r}(t_i)}{\De t_i}}\De t_i,
\end{equation}
que en el l\'imite es la integral
\begin{equation}
  \label{eq:longitud_arco}
  S(t)=\int_{t_0}^{t}\norm{\dv{}{t'}\vb{r}(t')}\dd t' .
\end{equation}

Por el teorema fundamental del c\'alculo,
\begin{equation}
  \label{eq:derivada_arco}
  \dv{}{t}S(t)=\norm{\dv{}{t}\vb{r}(t)}>0
  \qquad\text{para}\quad t>t_0,
\end{equation}
de modo que $S$ es mon\'otonamente creciente. Vale entonces el teorema de la
funci\'on inversa y podemos escribir $t=t(s)$, es decir,
\textbf{parametrizar por longitud de arco}. Derivando la identidad
$t(s(t))=t$ por la regla de la cadena,
\begin{equation}
  \label{eq:inversa_arco}
  \dv{}{s}t=\frac{1}{\dv{}{t}s} .
\end{equation}

\begin{ejemplo}[Circunferencia]
  \label{ej:arco_circulo}
  Para $\vb{r}(\vph)=(R\cos\vph,\,R\sin\vph)^\tp$ con
  $\vph\in[0,\vph_1]$,
  \begin{equation}
    \label{eq:arco_circulo}
    \begin{aligned}
      S&=\int_{0}^{\vph_1}\norm{\dv{}{\vph}\vb{r}}\dd\vph\\
       &=\int_{0}^{\vph_1}\sqrt{\qty{\dv{x}{\vph}}^2
          +\qty{\dv{y}{\vph}}^2}\dd\vph=R\vph_1 .
    \end{aligned}
  \end{equation}
\end{ejemplo}

\begin{ejemplo}[H\'elice]
  \label{ej:arco_helice}
  Para \eqref{eq:helice},
  \begin{equation}
    \label{eq:arco_helice}
    \begin{aligned}
      \dv{}{t}\vb{r}(t)&=
      \begin{pmatrix}
        -R\om\sin(\om t)\\
        R\om\cos(\om t)\\
        a
      \end{pmatrix},\\
      S(t)&=\int_{0}^{t}\sqrt{R^2\om^2+a^2}\dd t'
           =\sqrt{R^2\om^2+a^2}\;t,
    \end{aligned}
  \end{equation}
  Escribamos $\La=\sqrt{R^2\om^2+a^2}$ para abreviar; entonces
  \begin{equation}
    \label{eq:t_de_s_helice}
    t=\frac{s}{\La},
  \end{equation}
  que describe la din\'amica de la curva. Claramente
  \begin{equation}
    \label{eq:derivadas_arco_helice}
    \dv{}{t}s=\La>0
    \imp
    \dv{}{s}t=\frac{1}{\La}>0 .
  \end{equation}
  Sustituyendo \eqref{eq:t_de_s_helice} en \eqref{eq:helice}, y con
  $\vph(s)=s\om/\La$, se obtiene la h\'elice parametrizada por longitud de
  arco,
  \begin{equation}
    \label{eq:helice_arco}
    \vb{r}(s)=
    \begin{pmatrix}
      R\cos\vph(s)\\
      R\sin\vph(s)\\
      a s/\La
    \end{pmatrix}.
  \end{equation}
\end{ejemplo}

\subsection{La triada de Fr\'enet--Serret}
\label{subsec:frenet_serret}

La parametrizaci\'on por longitud de arco permite construir una teor\'ia
robusta de curvas. Por la regla de la cadena,
\begin{equation}
  \label{eq:tangente_unitario}
  \dv{}{s}\vb{r}(s)=\dv{}{t}\vb{r}\,\dv{}{s}t
    =\frac{\vb{V}}{\norm{\vb{V}}}=\vb{T}(s),
\end{equation}
el \textbf{vector tangente unitario}, que cumple $\norm{\vb{T}}=1$.

De ah\'i se sigue una consecuencia inmediata: como
$\vb{T}\cdot\vb{T}=1$ es constante,
\begin{equation}
  \label{eq:T_perp_Tprima}
  \dv{}{s}\qty{\vb{T}\cdot\vb{T}}
    =\vb{T}'\cdot\vb{T}+\vb{T}\cdot\vb{T}'
    =2\,\vb{T}'\cdot\vb{T}=0
  \imp \vb{T}'\perp\vb{T},
\end{equation}
donde la prima denota derivada respecto de $s$.

\begin{definicion}[Normal y curvatura]
  \label{def:normal_curvatura}
  El \textbf{vector normal} $\vb{N}$ es el unitario perpendicular a $\vb{T}$
  definido por
  \begin{equation}
    \label{eq:normal}
    \vb{T}'=\ka(s)\vb{N},
    \qquad
    \ka(s)=\norm{\vb{T}'},
  \end{equation}
  y $\ka(s)$ es la \textbf{curvatura} de la curva.
\end{definicion}

\begin{definicion}[Binormal]
  \label{def:binormal}
  El \textbf{vector binormal} se define como
  \begin{equation}
    \label{eq:binormal}
    \vb{B}=\vb{T}\wedge\vb{N}.
  \end{equation}
  La terna $\set{\vb{T},\vb{N},\vb{B}}$ es el \textbf{marco de
  Fr\'enet--Serret}, local en cada punto de la curva.
\end{definicion}

Los tres vectores son unitarios,
$\norm{\vb{T}}=\norm{\vb{N}}=\norm{\vb{B}}=1$, y mutuamente ortogonales.

\subsection{Torsi\'on}
\label{subsec:torsion}

De $\vb{B}\cdot\vb{B}=1$ se obtiene, igual que en \eqref{eq:T_perp_Tprima},
\begin{equation}
  \label{eq:B_perp_Bprima}
  \dv{}{s}\qty{\vb{B}\cdot\vb{B}}=0
  \imp \vb{B}'\cdot\vb{B}=0
  \imp \vb{B}'\perp\vb{B}.
\end{equation}
Por otro lado, $\vb{B}\cdot\vb{T}=0$, de modo que
\begin{equation}
  \label{eq:B_perp_T}
  \begin{aligned}
    \dv{}{s}\qty{\vb{B}\cdot\vb{T}}&=0\\
    \imp \vb{B}'\cdot\vb{T}+\vb{B}\cdot\vb{T}'&=0 .
  \end{aligned}
\end{equation}
Pero $\vb{T}'=\ka(s)\vb{N}$ y $\vb{N}\perp\vb{B}$, luego el segundo t\'ermino
se anula y $\vb{B}'\cdot\vb{T}=0$, es decir $\vb{B}'\perp\vb{T}$. Siendo
$\vb{B}'$ perpendicular tanto a $\vb{B}$ como a $\vb{T}$, debe ser paralelo a
$\vb{N}$.

\begin{definicion}[Torsi\'on]
  \label{def:torsion}
  La \textbf{torsi\'on} $\ta(s)$ se define por
  \begin{equation}
    \label{eq:torsion}
    \vb{B}'=\ta(s)\vb{N},
    \qquad
    \abs{\ta(s)}=\norm{\vb{B}'} .
  \end{equation}
\end{definicion}

Falta la derivada de $\vb{N}$. Como $\vb{N}'\perp\vb{N}$, vive en el plano
generado por $\vb{T}$ y $\vb{B}$, es decir
$\vb{N}'=N_1\vb{T}+N_2\vb{B}$. Los coeficientes se obtienen derivando las
relaciones de ortogonalidad:
\begin{equation}
  \label{eq:coeficientes_Nprima}
  \begin{aligned}
    \vb{N}\cdot\vb{T}=0
      &\imp \vb{N}'\cdot\vb{T}+\vb{N}\cdot\vb{T}'=0\\
      &\imp \vb{N}'\cdot\vb{T}=-\ka(s),\\
    \vb{N}\cdot\vb{B}=0
      &\imp \vb{N}'\cdot\vb{B}+\vb{N}\cdot\vb{B}'=0\\
      &\imp \vb{N}'\cdot\vb{B}=-\ta(s).
  \end{aligned}
\end{equation}

\subsection{Recetario}
\label{subsec:recetario_frenet}

Reuniendo \eqref{eq:normal}, \eqref{eq:coeficientes_Nprima} y
\eqref{eq:torsion}, las \textbf{ecuaciones de Fr\'enet--Serret} son
\begin{equation}
  \label{eq:frenet_serret}
  \left\{
  \begin{aligned}
    \vb{T}'&=\ka(s)\vb{N},\\
    \vb{N}'&=-\ka(s)\vb{T}-\ta(s)\vb{B},\\
    \vb{B}'&=\ta(s)\vb{N}.
  \end{aligned}
  \right.
\end{equation}

Por el teorema de la funci\'on inversa se puede escribir la triada en
t\'erminos del par\'ametro original $t$, y entonces
\begin{align}
  \label{eq:curvatura_t}
  \ka(t)&=\frac{\norm{\dot{\vb{r}}(t)\wedge\ddot{\vb{r}}(t)}}
               {\norm{\dot{\vb{r}}(t)}^{3}},\\
  \label{eq:torsion_t}
  \ta(t)&=\frac{\dot{\vb{r}}(t)\cdot
                \qty{\ddot{\vb{r}}(t)\wedge\dddot{\vb{r}}(t)}}
               {\norm{\dot{\vb{r}}(t)\wedge\ddot{\vb{r}}(t)}^{2}},
\end{align}
donde
\begin{equation}
  \label{eq:puntos_derivadas}
  \dot{\vb{r}}=\dv{}{t}\vb{r}(t),
  \quad
  \ddot{\vb{r}}=\ddv{}{t}\vb{r}(t),
  \quad
  \dddot{\vb{r}}=\dnv{\vb{r}}{t}{3}(t).
\end{equation}

\begin{ejemplo}[Triada de la h\'elice]
  \label{ej:triada_helice}
  Con la notaci\'on de \eqref{eq:helice_arco}, derivando,
  \begin{equation}
    \label{eq:T_helice}
    \vb{T}=\dv{}{s}\vb{r}(s)=\frac{1}{\La}
    \begin{pmatrix}
      -R\om\sin\vph(s)\\
      R\om\cos\vph(s)\\
      a
    \end{pmatrix}.
  \end{equation}
  Derivando otra vez,
  \begin{equation}
    \label{eq:Tprima_helice}
    \vb{T}'=\underbrace{\frac{R\om^2}{R^2\om^2+a^2}}_{\ka(s)}
    \underbrace{
    \begin{pmatrix}
      -\cos\vph(s)\\
      -\sin\vph(s)\\
      0
    \end{pmatrix}}_{\vb{N}},
  \end{equation}
  de modo que la curvatura de la h\'elice es constante,
  $\ka=R\om^2/(R^2\om^2+a^2)$. El binormal resulta
  \begin{equation}
    \label{eq:B_helice}
    \vb{B}=\vb{T}\wedge\vb{N}=\frac{1}{\La}
    \begin{pmatrix}
      a\sin\vph(s)\\
      -a\cos\vph(s)\\
      R\om
    \end{pmatrix},
  \end{equation}
  y derivando,
  \begin{equation}
    \label{eq:Bprima_helice}
    \vb{B}'=\frac{a\om}{R^2\om^2+a^2}
    \begin{pmatrix}
      \cos\vph(s)\\
      \sin\vph(s)\\
      0
    \end{pmatrix}
    =-\frac{a\om}{R^2\om^2+a^2}\,\vb{N},
  \end{equation}
  es decir, la torsi\'on tambi\'en es constante,
  $\ta=-a\om/(R^2\om^2+a^2)$.
\end{ejemplo}
